%% kapa: exam=1 round=37 year=2026 period=1 form=A subject=civil_law subject_ko=민법 seq=23 number=23
%% source: Q-Net 감정평가사 기출문제 (cst00309) artlSeq=5258924 file=raw/1008/5258924/2026년 제37회 감정평가사 1차 1교시.pdf
%% conversion: kordoc 4.24.51 -> tools/gen1.py (auto); review=proofread-sonnet; flags=none
\begin{question}{23}
소유권이전등기청구권을 보전하기 위한 가등기를 한 경우에 관한 설명으로 옳지 않은 것은? (다툼이 있으면 판례에 따름)
\begin{choices}
  \item 가등기된 권리의 이전등기는 가등기에 대한 부기등기의 형식으로 할 수 있다.
  \item 무효인 가등기의 유용은 등기상 이해관계인이 없는 한 허용된다.
  \item 소유권이전등기청구권 보전을 위한 가등기가 있다고 하여, 소유권이전등기를 청구할 어떤 법률관계가 있다고 추정되지 않는다.
  \item 가등기권리자는 무효인 중복소유권보존등기의 명의인을 상대로 그 보존등기의 말소를 구할 수 있다.
  \item 가등기에 기하여 본등기가 경료되더라도 본등기에 의한 물권변동의 효력이 가등기한 때로 소급하지 않는다.
\end{choices}
\end{question}
\answer{4}
